% Assignment source: ne630/homework/html/HW08.html.
% Legacy foundation: ne630_problems/lesson_08/statement.tex and
% ne630_problems/lesson_08/HW08_solution.ipynb.
% Problem 2 keeps the legacy notebook's computed eta values and adds the
% Fall 2026 interpretation prompt.

\noindent{\bf Problem 1}. Following the lecture example, use data from OpenMC
to reproduce Figure 3.2 from FNRP, and add a curve for U enriched to 4\%.
This figure should be produced using computer software (e.g., Python, Excel,
MATLAB) and {\it not} by hand.

\begin{adaptedsolution}
A correct solution computes
\[
 \eta(E)=
 \frac{e\,\bar{\nu}_{235}(E)\sigma_{f,235}(E)
       +(1-e)\,\bar{\nu}_{238}(E)\sigma_{f,238}(E)}
      {e\,\sigma_{a,235}(E)+(1-e)\,\sigma_{a,238}(E)}
\]
for the requested enrichments, where $e$ is the atom fraction of
\isoA{U}{235}.  The absorption cross sections should include capture plus
fission:
\[
 \sigma_a(E)=\sigma_\gamma(E)+\sigma_f(E).
\]

The plotted curves should reproduce the main behavior of FNRP Fig.~3.2: the
natural-uranium curve remains below the enriched cases over most energies, the
20\% curve is highest, and the 4\% curve lies between them.  The raw evaluated
data are noisy where the pointwise cross sections and $\bar{\nu}$ are tabulated
on different energy grids, so interpolation to a common grid is expected.
Smoothing the plotted curve is acceptable for visual comparison, but grading
should be based on the unsmoothed values used in any numerical follow-up.
\end{adaptedsolution}

\newpage

\noindent{\bf Problem 2}. For the 4\% values, compute the expected value of
$\eta(E)$ using
\[
p(E)=
  \begin{cases}
    C, & E_{\mathrm{min}}\le E\le E_{\mathrm{max}},\\
    0, & \text{otherwise},
  \end{cases}
\]
for the following three cases:
\begin{enumerate}[label=(\alph*)]
\item $E_{\mathrm{min}}=10^{-3}~\mathrm{eV}$ and
      $E_{\mathrm{max}}=1~\mathrm{eV}$ (the ``thermal'' range).
\item $E_{\mathrm{min}}=1~\mathrm{eV}$ and
      $E_{\mathrm{max}}=10^5~\mathrm{eV}$ (the ``epithermal'' range).
\item $E_{\mathrm{min}}=10^5~\mathrm{eV}$ and
      $E_{\mathrm{max}}=10^7~\mathrm{eV}$ (the ``fast'' range).
\end{enumerate}
For each case, define $C$ using the properties of probability density
functions; that is,
\[
 \int_{E_{\mathrm{min}}}^{E_{\mathrm{max}}}p(E)\,dE=1.
\]
Generally, reactors are classified as being either ``thermal'' or ``fast''
spectrum reactors. Based on your values for $\bar{\eta}$ in each energy range,
explain why we do not often consider ``intermediate'' spectrum reactors.

\begin{adaptedsolution}
For a uniform probability density over the interval
$[E_{\min},E_{\max}]$,
\[
 C=\frac{1}{E_{\max}-E_{\min}},
 \qquad
 \bar{\eta}
 =\frac{1}{E_{\max}-E_{\min}}
  \int_{E_{\min}}^{E_{\max}}\eta(E)\,dE.
\]
Numerically, the legacy solution notebook evaluates $\eta(E)$ from pointwise
BNL/OpenMC-style data and integrates with the trapezoid rule.  For the full
range $10^{-3}\le E\le10^7~\mathrm{eV}$, the same calculation gives
\[
 \bar{\eta}_{\mathrm{uniform}}\approx 2.82.
\]
If instead the flux-weighting density is taken proportional to $1/E$ over
the full range, the notebook gives
\[
 \bar{\eta}_{1/E}
 =\frac{\int \eta(E)E^{-1}\,dE}{\int E^{-1}\,dE}
 \approx 1.36.
\]

Using the local ENDF/B-VIII.1 OpenMC HDF5 data and a converged logarithmic
energy grid, the three assigned uniform-range averages are approximately
\[
\begin{array}{ll}
10^{-3}\le E\le1~\mathrm{eV}: & \bar{\eta}\approx1.82,\\
1\le E\le10^5~\mathrm{eV}: & \bar{\eta}\approx0.552,\\
10^5\le E\le10^7~\mathrm{eV}: & \bar{\eta}\approx2.85.
\end{array}
\]
The epithermal/intermediate range is dominated by strong resonance absorption
structure and gives a much lower average reproduction factor than either the
thermal or fast range in this calculation.  That is the main physical reason
reactor spectra are usually pushed toward thermal or fast designs rather than
left in an intermediate spectrum.
\end{adaptedsolution}

\newpage

\noindent{\bf Problem 3}. The slowing-down properties of three common
moderators are provided in Table 3.1 of FNRP. Add a row to that table for Be,
another element with small mass, and compare Be as a moderator with
$\mathrm{H}_2\mathrm{O}$, $\mathrm{D}_2\mathrm{O}$, and graphite. You may use
data from Appendix E, Table E.3.

\begin{adaptedsolution}
Natural Be consists only of \isoA{Be}{9} and, from Table~E.3, has a mass
density of $\rho=1.85~\gram\per\centi\meter\cubed$.  Thus
\[
 \alpha=\left(\frac{A-1}{A+1}\right)^2
       =\left(\frac{8}{10}\right)^2=0.64,
\]
and
\[
 \boxed{\xi
 =1+\frac{\alpha}{1-\alpha}\ln\alpha
 \approx0.207}.
\]
The number density of natural Be is
\[
 n_{\mathrm{Be}}
 =\frac{1.85(6.022\times10^{23})}{9.013}
 \approx1.24\times10^{23}~\mathrm{atoms/cm^3}.
\]
Using $\sigma_s=6.151~\mathrm b$ and $\sigma_a=0.0085~\mathrm b$,
\[
 \Sigma_s\approx0.76~\centi\meter^{-1},
 \qquad
 \Sigma_a\approx1.1\times10^{-3}~\centi\meter^{-1}.
\]
Therefore,
\[
 \boxed{\xi\Sigma_s\approx0.157~\centi\meter^{-1}},
 \qquad
 \boxed{\frac{\xi\Sigma_s}{\Sigma_a}\approx1.5\times10^2}.
\]
Be has a slowing-down decrement and slowing-down power between heavy water and
graphite, as expected from its mass.  Its slowing-down ratio is better than
light water but below graphite.  It is plausible as a reflector or specialty
moderating material, but graphite and heavy water remain the more familiar
natural-uranium moderator choices.
\end{adaptedsolution}

\clearpage
